\documentclass[12pt,letterpaper]{article}

\newenvironment{proof}{\noindent{\bf Proof:}}{\qed\bigskip}

\newtheorem{theorem}{Theorem}
\newtheorem{corollary}{Corollary}
\newtheorem{lemma}{Lemma} 
\newtheorem{claim}{Claim}
\newtheorem{fact}{Fact}
\newtheorem{definition}{Definition}
\newtheorem{assumption}{Assumption}
\newtheorem{observation}{Observation}
\newtheorem{example}{Example}
\newcommand{\qed}{\rule{7pt}{7pt}}

\newcommand{\lecture}[5]{
	\thispagestyle{plain} 
	\newpage
	\setcounter{page}{1}
	\noindent
	\begin{center}
		\framebox{ \vbox{ \hbox to 6.28in
				{\bf AI543: Foundations of Online Machine Learning  \hfill #1}
				\vspace{4mm}
				\hbox to 6.28in
				{\hspace{2.5in}\large\mbox{Lecture #2}}
				\vspace{4mm}
				\hbox to 6.28in
				{{\it Instructor: #3 \hfill Date: #4 \hfill Scriber: #5}}
			}}
		\end{center}
	}

\newcommand{\assignment}[4]{
\thispagestyle{plain} 
\newpage
\setcounter{page}{1}
\noindent
\begin{center}
\framebox{ \vbox{ \hbox to 6.28in
{\bf AI543: Foundations of Online Machine Learning   \hfill #1}
\vspace{4mm}
\hbox to 6.28in
{\hspace{2.5in}\large\mbox{Homework #2}}
\vspace{4mm}
\hbox to 6.28in
{{\it Handed Out: #3 \hfill Due: #4}}
}}
\end{center}
}

\newcommand{\homework}[3]{
	\thispagestyle{plain} 
	\newpage
	\setcounter{page}{1}
	\noindent
	\begin{center}
		\framebox{ \vbox{ \hbox to 6.28in
				{\bf AI543: Foundations of Online Machine Learning   \hfill #1}
				\vspace{4mm}
				\hbox to 6.28in
				{\hspace{2.5in}\large\mbox{Homework #2}}
				\vspace{4mm}
				\hbox to 6.28in
				{\it OSU ID: 934--632--976 \hfill Student Name: Mrinal Bharadwaj}
                \hbox to 6.28in
                {\it Date: January 28, 2026 \hfill}
			}}
		\end{center}
	}

\newcommand{\solution}[4]{
\thispagestyle{plain} 
\newpage
\setcounter{page}{1}
\noindent
\begin{center}
\framebox{ \vbox{ \hbox to 6.28in
{\bf AI543: Foundations of Online Machine Learning   \hfill #4}
\vspace{4mm}
\hbox to 6.28in
{\hspace{2.5in}\large\mbox{Homework Solution #3}}
\vspace{4mm}
\hbox to 6.28in
{#1 \hfill {\it Handed In: #2}}
}}
\end{center}
\markright{#1}
}

\newenvironment{algorithm}
{\begin{center}
\begin{tabular}{|l|}
\hline
\begin{minipage}{1in}
\begin{tabbing}
\quad\=\qquad\=\qquad\=\qquad\=\qquad\=\qquad\=\qquad\=\kill}
{\end{tabbing}
\end{minipage} \\
\hline
\end{tabular}
\end{center}}

\def\Comment#1{\textsf{\textsl{$\langle\!\langle$#1\/$\rangle\!\rangle$}}}


\usepackage{graphicx,amssymb,amsmath}
\sloppy
\newcommand{\ignore}[1]{}

\oddsidemargin 0in
\evensidemargin 0in
\textwidth 6.5in
\topmargin -0.5in
\textheight 9.0in

\begin{document}

\homework{Winter 2026}{$1$}{XXX} %Replace XXX with your name


\begin{enumerate}
	
	\item[1.] 
    \textbf{Problem:} Let $C$ be the solution set of a quadratic inequality:
\[
C = \{\mathbf{x} \in \mathbb{R}^d : \mathbf{x}^\top A\mathbf{x} + \mathbf{b}^\top \mathbf{x} + c \leq 0\},
\]
where $\mathbf{b} \in \mathbb{R}^d$, $c \in \mathbb{R}$, $A \in \mathbb{R}^{d \times d}$ is a symmetric matrix. Show that $C$ is convex if $A \succ 0$.

\bigskip

\textbf{Solution:}

\begin{proof}
Define the function $f: \mathbb{R}^d \to \mathbb{R}$ by
\[
f(\mathbf{x}) = \mathbf{x}^\top A\mathbf{x} + \mathbf{b}^\top \mathbf{x} + c.
\]

Note that $C = \{\mathbf{x} \in \mathbb{R}^d : f(\mathbf{x}) \leq 0\}$ is the sublevel set of $f$ at level 0.

\textbf{Step 1: Show that $f$ is convex when $A \succ 0$.}

We compute the Hessian of $f$. The gradient is:
\[
\nabla f(\mathbf{x}) = 2A\mathbf{x} + \mathbf{b}.
\]

The Hessian is:
\[
\nabla^2 f(\mathbf{x}) = 2A.
\]

Since $A \succ 0$ (positive definite), we have $2A \succ 0$, which means the Hessian is positive definite for all $\mathbf{x}$.

Therefore, $f$ is a strictly convex function.

\textbf{Step 2: Sublevel sets of convex functions are convex.}

Let $\mathbf{x}, \mathbf{y} \in C$ and $\lambda \in [0, 1]$. We need to show that $\lambda \mathbf{x} + (1-\lambda)\mathbf{y} \in C$.

Since $\mathbf{x}, \mathbf{y} \in C$, we have:
\[
f(\mathbf{x}) \leq 0 \quad \text{and} \quad f(\mathbf{y}) \leq 0.
\]

Since $f$ is convex, by the definition of convexity:
\[
f(\lambda \mathbf{x} + (1-\lambda)\mathbf{y}) \leq \lambda f(\mathbf{x}) + (1-\lambda)f(\mathbf{y}).
\]

Using our bounds:
\[
f(\lambda \mathbf{x} + (1-\lambda)\mathbf{y}) \leq \lambda f(\mathbf{x}) + (1-\lambda)f(\mathbf{y}) \leq \lambda \cdot 0 + (1-\lambda) \cdot 0 = 0.
\]

Therefore, $\lambda \mathbf{x} + (1-\lambda)\mathbf{y} \in C$.

\textbf{Conclusion:} Since for any two points in $C$ and any $\lambda \in [0,1]$, their convex combination also lies in $C$, the set $C$ is convex.

\end{proof}

	
	\item[2.] 
	\textbf{Problem:} Show that
\[
f(x) = \frac{1}{2}\|x\|_2^2, \quad x \in \mathbb{R}^d,
\]
is 1-strongly convex w.r.t.\ $\|\cdot\|_2$.

\bigskip

\textbf{Solution:}

\begin{proof}
We will provide three different approaches to prove this result.

\bigskip

\textbf{Method 1: Using the Hessian (Second-Order Condition)}

A twice-differentiable function $f$ is $\mu$-strongly convex if and only if 
\[
\nabla^2 f(x) \succeq \mu I \quad \text{for all } x.
\]

For $f(x) = \frac{1}{2}\|x\|_2^2 = \frac{1}{2}x^\top x$:

The gradient is:
\[
\nabla f(x) = x.
\]

The Hessian is:
\[
\nabla^2 f(x) = I,
\]
where $I$ is the $d \times d$ identity matrix.

Since $\nabla^2 f(x) = I \succeq 1 \cdot I$, the function $f$ is 1-strongly convex.

\bigskip

\textbf{Method 2: Using the Definition of Strong Convexity}

A function $f$ is $\mu$-strongly convex w.r.t.\ $\|\cdot\|_2$ if for all $x, y \in \mathbb{R}^d$ and $\lambda \in [0,1]$:
\[
f(\lambda x + (1-\lambda)y) \leq \lambda f(x) + (1-\lambda)f(y) - \frac{\mu}{2}\lambda(1-\lambda)\|x - y\|_2^2.
\]

We need to verify this for $\mu = 1$. Let $z = \lambda x + (1-\lambda)y$. Then:
\[
f(z) = \frac{1}{2}\|\lambda x + (1-\lambda)y\|_2^2.
\]

Expanding the right-hand side:
\begin{align*}
\frac{1}{2}\|\lambda x + (1-\lambda)y\|_2^2 &= \frac{1}{2}\left[\lambda^2\|x\|_2^2 + 2\lambda(1-\lambda)x^\top y + (1-\lambda)^2\|y\|_2^2\right].
\end{align*}

Now, compute $\lambda f(x) + (1-\lambda)f(y) - \frac{1}{2}\lambda(1-\lambda)\|x-y\|_2^2$:
\begin{align*}
&= \frac{\lambda}{2}\|x\|_2^2 + \frac{1-\lambda}{2}\|y\|_2^2 - \frac{\lambda(1-\lambda)}{2}\|x-y\|_2^2 \\
&= \frac{\lambda}{2}\|x\|_2^2 + \frac{1-\lambda}{2}\|y\|_2^2 - \frac{\lambda(1-\lambda)}{2}\left(\|x\|_2^2 - 2x^\top y + \|y\|_2^2\right) \\
&= \frac{1}{2}\left[\lambda\|x\|_2^2 + (1-\lambda)\|y\|_2^2 - \lambda(1-\lambda)\|x\|_2^2 + 2\lambda(1-\lambda)x^\top y - \lambda(1-\lambda)\|y\|_2^2\right] \\
&= \frac{1}{2}\left[\lambda(1-(1-\lambda))\|x\|_2^2 + (1-\lambda)(1-\lambda)\|y\|_2^2 + 2\lambda(1-\lambda)x^\top y\right] \\
&= \frac{1}{2}\left[\lambda^2\|x\|_2^2 + (1-\lambda)^2\|y\|_2^2 + 2\lambda(1-\lambda)x^\top y\right] \\
&= \frac{1}{2}\|\lambda x + (1-\lambda)y\|_2^2 = f(z).
\end{align*}

The inequality holds with equality, confirming that $f$ is 1-strongly convex.

\bigskip

\textbf{Method 3: Using the First-Order Condition}

A differentiable function $f$ is $\mu$-strongly convex if and only if for all $x, y \in \mathbb{R}^d$:
\[
f(y) \geq f(x) + \nabla f(x)^\top (y - x) + \frac{\mu}{2}\|y - x\|_2^2.
\]

With $\nabla f(x) = x$, we need to show:
\[
\frac{1}{2}\|y\|_2^2 \geq \frac{1}{2}\|x\|_2^2 + x^\top(y - x) + \frac{1}{2}\|y - x\|_2^2.
\]

Expanding the right-hand side:
\begin{align*}
\text{RHS} &= \frac{1}{2}\|x\|_2^2 + x^\top y - \|x\|_2^2 + \frac{1}{2}\left(\|y\|_2^2 - 2x^\top y + \|x\|_2^2\right) \\
&= \frac{1}{2}\|x\|_2^2 + x^\top y - \|x\|_2^2 + \frac{1}{2}\|y\|_2^2 - x^\top y + \frac{1}{2}\|x\|_2^2 \\
&= \frac{1}{2}\|y\|_2^2 = \text{LHS}.
\end{align*}

The inequality holds with equality, confirming 1-strong convexity.

\bigskip

\textbf{Conclusion:} By all three methods, we have shown that $f(x) = \frac{1}{2}\|x\|_2^2$ is 1-strongly convex with respect to $\|\cdot\|_2$.

\end{proof}
	\item[3.]
    \textbf{Problem:} Recall the definition of $p$-norms as
\[
\|x\|_p = \left(\sum_{i=1}^{d}|x_i|^p\right)^{1/p}
\]
for $p \geq 1$.

\begin{enumerate}
    \item Show that 2-norms are always convex.
    \item Show that $p$-norms are always convex. [Hint: Apply Minkowski inequality]
\end{enumerate}

\bigskip

\textbf{Solution:}

\bigskip

\textbf{Part (1): Show that 2-norms are always convex.}

\begin{proof}
A function $f: \mathbb{R}^d \to \mathbb{R}$ is convex if for all $x, y \in \mathbb{R}^d$ and $\lambda \in [0,1]$:
\[
f(\lambda x + (1-\lambda)y) \leq \lambda f(x) + (1-\lambda)f(y).
\]

Let $f(x) = \|x\|_2 = \sqrt{\sum_{i=1}^{d} x_i^2}$. We need to show:
\[
\|\lambda x + (1-\lambda)y\|_2 \leq \lambda\|x\|_2 + (1-\lambda)\|y\|_2.
\]

\textbf{Method 1: Using the Triangle Inequality}

The 2-norm satisfies the triangle inequality:
\[
\|a + b\|_2 \leq \|a\|_2 + \|b\|_2 \quad \text{for all } a, b \in \mathbb{R}^d.
\]

Let $a = \lambda x$ and $b = (1-\lambda)y$. Then:
\begin{align*}
\|\lambda x + (1-\lambda)y\|_2 &\leq \|\lambda x\|_2 + \|(1-\lambda)y\|_2 \\
&= |\lambda| \cdot \|x\|_2 + |1-\lambda| \cdot \|y\|_2 \\
&= \lambda \|x\|_2 + (1-\lambda)\|y\|_2,
\end{align*}
where the last equality holds because $\lambda \in [0,1]$, so $\lambda \geq 0$ and $1-\lambda \geq 0$.

This proves that $\|x\|_2$ is convex.

\bigskip

\textbf{Method 2: Using the Cauchy-Schwarz Inequality}

The triangle inequality for the 2-norm can be derived from the Cauchy-Schwarz inequality. For any $a, b \in \mathbb{R}^d$:
\[
\|a + b\|_2^2 = \|a\|_2^2 + 2a^\top b + \|b\|_2^2 \leq \|a\|_2^2 + 2\|a\|_2\|b\|_2 + \|b\|_2^2 = (\|a\|_2 + \|b\|_2)^2,
\]
where we used $a^\top b \leq \|a\|_2\|b\|_2$ (Cauchy-Schwarz).

Taking square roots: $\|a + b\|_2 \leq \|a\|_2 + \|b\|_2$.

The convexity then follows as in Method 1.
\end{proof}

\bigskip

\textbf{Part (2): Show that $p$-norms are always convex for $p \geq 1$.}

\begin{proof}
We use the \textbf{Minkowski Inequality}, which states that for $p \geq 1$ and any $a, b \in \mathbb{R}^d$:
\[
\|a + b\|_p \leq \|a\|_p + \|b\|_p.
\]

This is the triangle inequality for the $p$-norm.

\bigskip

To prove convexity of $f(x) = \|x\|_p$, let $x, y \in \mathbb{R}^d$ and $\lambda \in [0,1]$.

Set $a = \lambda x$ and $b = (1-\lambda)y$. Applying the Minkowski inequality:
\begin{align*}
\|\lambda x + (1-\lambda)y\|_p &\leq \|\lambda x\|_p + \|(1-\lambda)y\|_p \\
&= |\lambda| \cdot \|x\|_p + |1-\lambda| \cdot \|y\|_p \\
&= \lambda \|x\|_p + (1-\lambda)\|y\|_p,
\end{align*}
where we used the absolute homogeneity of norms: $\|\alpha z\|_p = |\alpha| \cdot \|z\|_p$.

Since $\lambda \in [0,1]$, both $\lambda \geq 0$ and $1-\lambda \geq 0$, so:
\[
\|\lambda x + (1-\lambda)y\|_p \leq \lambda \|x\|_p + (1-\lambda)\|y\|_p.
\]

This is precisely the definition of convexity.

\bigskip

\textbf{Remark (Proof of Minkowski Inequality):}

For completeness, the Minkowski inequality is typically proved using H\"older's inequality. For $p > 1$, let $q$ be the conjugate exponent satisfying $\frac{1}{p} + \frac{1}{q} = 1$.

H\"older's inequality states:
\[
\sum_{i=1}^{d}|a_i b_i| \leq \|a\|_p \|b\|_q.
\]

Using this, one can show:
\begin{align*}
\|a + b\|_p^p &= \sum_{i=1}^{d}|a_i + b_i|^p \\
&= \sum_{i=1}^{d}|a_i + b_i|^{p-1}|a_i| + \sum_{i=1}^{d}|a_i + b_i|^{p-1}|b_i| \\
&\leq \||a+b|^{p-1}\|_q \cdot \|a\|_p + \||a+b|^{p-1}\|_q \cdot \|b\|_p \\
&= \|a + b\|_p^{p/q}(\|a\|_p + \|b\|_p).
\end{align*}

Since $p - \frac{p}{q} = 1$, dividing both sides by $\|a+b\|_p^{p/q}$ yields:
\[
\|a + b\|_p \leq \|a\|_p + \|b\|_p.
\]

\bigskip

\textbf{Conclusion:} For all $p \geq 1$, the $p$-norm $\|x\|_p$ is a convex function, as a direct consequence of the Minkowski inequality (triangle inequality for $p$-norms).

\end{proof}
	
	\item[4.]
\textbf{Problem:} Recall the Online Gradient Descent algorithm. We know that setting a constant learning rate $\eta^* = \frac{D}{L\sqrt{T}}$ gives us regret $DL\sqrt{T}$. We can achieve similar regret with a dynamic learning rate.

\begin{enumerate}
    \item Prove that $\sum_{t=1}^{T}\frac{1}{\sqrt{t}} \leq 2\sqrt{T} - 1$
    \item Show the regret of OGD if we choose $\eta_t \propto \frac{1}{\sqrt{t}}$. [Hint: use the above inequality.]
\end{enumerate}

\bigskip

\textbf{Solution:}

\bigskip

\textbf{Part (1): Prove that $\displaystyle\sum_{t=1}^{T}\frac{1}{\sqrt{t}} \leq 2\sqrt{T} - 1$}

\begin{proof}
We use the integral comparison technique. Since $f(t) = \frac{1}{\sqrt{t}}$ is a decreasing function for $t > 0$, we have:
\[
\frac{1}{\sqrt{t}} \leq \int_{t-1}^{t} \frac{1}{\sqrt{x}}\, dx \quad \text{for } t \geq 2.
\]

This is because the function value at $t$ (the right endpoint) is less than or equal to the average value over the interval $[t-1, t]$ for a decreasing function.

Therefore:
\begin{align*}
\sum_{t=1}^{T}\frac{1}{\sqrt{t}} &= \frac{1}{\sqrt{1}} + \sum_{t=2}^{T}\frac{1}{\sqrt{t}} \\
&\leq 1 + \sum_{t=2}^{T}\int_{t-1}^{t}\frac{1}{\sqrt{x}}\, dx \\
&= 1 + \int_{1}^{T}\frac{1}{\sqrt{x}}\, dx \\
&= 1 + \left[2\sqrt{x}\right]_{1}^{T} \\
&= 1 + 2\sqrt{T} - 2 \\
&= 2\sqrt{T} - 1.
\end{align*}

Thus, we have proven:
\[
\boxed{\sum_{t=1}^{T}\frac{1}{\sqrt{t}} \leq 2\sqrt{T} - 1}
\]
\end{proof}

\bigskip

\textbf{Part (2): Show the regret of OGD with $\eta_t \propto \frac{1}{\sqrt{t}}$}

\begin{proof}
Recall the regret bound for Online Gradient Descent with time-varying learning rates. From the standard OGD analysis, we have the following bound for the regret:
\[
R_T = \sum_{t=1}^{T}f_t(x_t) - \sum_{t=1}^{T}f_t(x^*) \leq \sum_{t=1}^{T}\left(\frac{\|x_t - x^*\|^2}{2\eta_t} - \frac{\|x_{t+1} - x^*\|^2}{2\eta_t} + \frac{\eta_t}{2}\|g_t\|^2\right),
\]
where $g_t = \nabla f_t(x_t)$.

With time-varying learning rates, using a telescoping argument and the fact that $\eta_t$ is decreasing, we get:
\[
R_T \leq \frac{D^2}{2\eta_T} + \frac{1}{2}\sum_{t=1}^{T}\eta_t \|g_t\|^2,
\]
where $D = \max_{x,y \in \mathcal{C}}\|x - y\|$ is the diameter of the constraint set.

\bigskip

\textbf{Setting the learning rate:} Let $\eta_t = \frac{D}{L\sqrt{t}}$, where $L$ is the Lipschitz constant satisfying $\|g_t\| \leq L$ for all $t$.

\bigskip

\textbf{Bounding the first term:}
\[
\frac{D^2}{2\eta_T} = \frac{D^2}{2 \cdot \frac{D}{L\sqrt{T}}} = \frac{D^2 \cdot L\sqrt{T}}{2D} = \frac{DL\sqrt{T}}{2}.
\]

The distributive property is applied properly:
\begin{itemize}
    \item $\frac{DL}{2} \times 2\sqrt{T} = DL\sqrt{T}$
    \item $\frac{DL}{2} \times 1 = \frac{DL}{2}$
\end{itemize}

\textbf{Bounding the second term:}
\begin{align*}
\frac{1}{2}\sum_{t=1}^{T}\eta_t\|g_t\|^2 &\leq \frac{1}{2}\sum_{t=1}^{T}\eta_t L^2 \\
&= \frac{L^2}{2}\sum_{t=1}^{T}\frac{D}{L\sqrt{t}} \\
&= \frac{DL}{2}\sum_{t=1}^{T}\frac{1}{\sqrt{t}} \\
&\leq \frac{DL}{2}(2\sqrt{T} - 1) \quad \text{(using Part 1)} \\
&= DL\sqrt{T} - \frac{DL}{2}.\text{(by distributive property)}
\end{align*}

In my bound for the second term, I have:
$$\frac{DL}{2}\sum_{t=1}^{T}\frac{1}{\sqrt{t}} \leq \frac{DL}{2}(2\sqrt{T} - 1)$$

When I expand $(2\sqrt{T} - 1)$:
$$\frac{DL}{2}(2\sqrt{T} - 1) = \frac{DL}{2} \cdot 2\sqrt{T} - \frac{DL}{2} \cdot 1$$

The DL/2 comes from:
$$\frac{DL}{2} \times (-1) = -\frac{DL}{2}$$

\textbf{Total regret bound:}
\begin{align*}
R_T &\leq \frac{DL\sqrt{T}}{2} + DL\sqrt{T} - \frac{DL}{2} \\
&= \frac{3}{2}DL\sqrt{T} - \frac{DL}{2} \\
&\leq \frac{3}{2}DL\sqrt{T}.
\end{align*}

\bigskip

\textbf{Conclusion:} With the dynamic learning rate $\eta_t = \frac{D}{L\sqrt{t}}$, the regret of Online Gradient Descent is:
\[
\boxed{R_T \leq \frac{3}{2}DL\sqrt{T} = O(DL\sqrt{T})}
\]

This achieves the same $O(\sqrt{T})$ regret rate as the constant learning rate $\eta^* = \frac{D}{L\sqrt{T}}$, but importantly, the dynamic learning rate does \textbf{not require knowledge of $T$ in advance}. This makes it an \textbf{anytime} algorithm.

\bigskip

\textbf{Remark:} The constant factor $\frac{3}{2}$ is slightly worse than the constant factor of $1$ achieved by the optimal constant learning rate, but this is a small price to pay for not needing to know the time horizon $T$ ahead of time.

\end{proof}

        \item[5.]
\textbf{Problem:} Recall FTL for number-guessing game. Extend the previous algorithm and analysis to the vector case when the adversary selects a vector $y_t \in \mathbb{R}^d$, $\|y\|_2 \leq 1$, the algorithm guesses a vector $x_t \in \mathbb{R}^d$, and the loss function is $(x_t - y_t)^2$. Show an upper bound to the regret logarithmic in $T$ that is independent from $d$.

\bigskip

\textbf{Solution:}

\begin{proof}

\textbf{Step 1: Define the FTL Algorithm for the Vector Case}

The loss at round $t$ is:
\[
\ell_t(x) = \|x - y_t\|_2^2 = \sum_{i=1}^{d}(x_i - y_{t,i})^2.
\]

The cumulative loss up to time $t-1$ is:
\[
L_{t-1}(x) = \sum_{s=1}^{t-1}\|x - y_s\|_2^2.
\]

The FTL algorithm chooses:
\[
x_t = \arg\min_{x \in \mathbb{R}^d} L_{t-1}(x) = \arg\min_{x \in \mathbb{R}^d} \sum_{s=1}^{t-1}\|x - y_s\|_2^2.
\]

\bigskip

\textbf{Step 2: Compute the FTL Prediction}

Taking the gradient and setting it to zero:
\[
\nabla_x L_{t-1}(x) = \sum_{s=1}^{t-1} 2(x - y_s) = 2(t-1)x - 2\sum_{s=1}^{t-1}y_s = 0.
\]

Solving for $x$:
\[
x_t = \frac{1}{t-1}\sum_{s=1}^{t-1}y_s = \bar{y}_{t-1},
\]
where $\bar{y}_{t-1}$ is the average of all previous adversary vectors.

For $t = 1$, we initialize $x_1 = 0$ (or any fixed point).

\bigskip

\textbf{Step 3: Use the FTL Lemma (Be-The-Leader Bound)}

The standard FTL lemma states:
\[
\sum_{t=1}^{T}\ell_t(x_t) - \sum_{t=1}^{T}\ell_t(x^*) \leq \sum_{t=1}^{T}\left[\ell_t(x_t) - \ell_t(x_{t+1})\right],
\]
where $x^* = \arg\min_x \sum_{t=1}^{T}\ell_t(x) = \bar{y}_T$ is the best fixed action in hindsight.

\bigskip

\textbf{Step 4: Bound the Instantaneous Regret}

We need to bound $\ell_t(x_t) - \ell_t(x_{t+1})$ for each $t$.

Recall:
\begin{align*}
x_t &= \frac{1}{t-1}\sum_{s=1}^{t-1}y_s = \bar{y}_{t-1} \quad (\text{for } t \geq 2), \\
x_{t+1} &= \frac{1}{t}\sum_{s=1}^{t}y_s = \bar{y}_t.
\end{align*}

We can write the update relation:
\[
x_{t+1} = \bar{y}_t = \frac{(t-1)\bar{y}_{t-1} + y_t}{t} = \frac{t-1}{t}x_t + \frac{1}{t}y_t.
\]

Therefore:
\[
x_{t+1} - x_t = \frac{1}{t}(y_t - x_t).
\]

Now compute:
\begin{align*}
\ell_t(x_t) - \ell_t(x_{t+1}) &= \|x_t - y_t\|_2^2 - \|x_{t+1} - y_t\|_2^2.
\end{align*}

Using $x_{t+1} = x_t + \frac{1}{t}(y_t - x_t) = \frac{t-1}{t}x_t + \frac{1}{t}y_t$:
\[
x_{t+1} - y_t = \frac{t-1}{t}x_t + \frac{1}{t}y_t - y_t = \frac{t-1}{t}(x_t - y_t).
\]

Thus:
\[
\|x_{t+1} - y_t\|_2^2 = \left(\frac{t-1}{t}\right)^2\|x_t - y_t\|_2^2.
\]

Therefore:
\begin{align*}
\ell_t(x_t) - \ell_t(x_{t+1}) &= \|x_t - y_t\|_2^2 - \left(\frac{t-1}{t}\right)^2\|x_t - y_t\|_2^2 \\
&= \|x_t - y_t\|_2^2 \left[1 - \left(\frac{t-1}{t}\right)^2\right] \\
&= \|x_t - y_t\|_2^2 \cdot \frac{t^2 - (t-1)^2}{t^2} \\
&= \|x_t - y_t\|_2^2 \cdot \frac{2t - 1}{t^2}.
\end{align*}

\bigskip

\textbf{Step 5: Bound $\|x_t - y_t\|_2^2$}

Since $x_t = \bar{y}_{t-1}$ is an average of vectors with $\|y_s\|_2 \leq 1$, we have $\|x_t\|_2 \leq 1$.

Also, $\|y_t\|_2 \leq 1$.

By the triangle inequality:
\[
\|x_t - y_t\|_2 \leq \|x_t\|_2 + \|y_t\|_2 \leq 2.
\]

Therefore:
\[
\|x_t - y_t\|_2^2 \leq 4.
\]

\bigskip

\textbf{Step 6: Compute the Total Regret Bound}

\begin{align*}
R_T &\leq \sum_{t=1}^{T}\left[\ell_t(x_t) - \ell_t(x_{t+1})\right] \\
&\leq \sum_{t=1}^{T} 4 \cdot \frac{2t-1}{t^2} \\
&= 4\sum_{t=1}^{T}\frac{2t-1}{t^2} \\
&= 4\sum_{t=1}^{T}\left(\frac{2}{t} - \frac{1}{t^2}\right) \\
&= 8\sum_{t=1}^{T}\frac{1}{t} - 4\sum_{t=1}^{T}\frac{1}{t^2}.
\end{align*}

Using the bounds:
\begin{itemize}
    \item $\sum_{t=1}^{T}\frac{1}{t} \leq 1 + \ln T$
    \item $\sum_{t=1}^{T}\frac{1}{t^2} \leq \frac{\pi^2}{6} < 2$ (converges)
\end{itemize}

We get:
\begin{align*}
R_T &\leq 8(1 + \ln T) - 4 \cdot 0 \\
&= 8 + 8\ln T.
\end{align*}

More precisely, since $\sum_{t=1}^{T}\frac{1}{t^2} > 0$:
\[
R_T \leq 8(1 + \ln T) = 8 + 8\ln T.
\]

\bigskip

\textbf{Conclusion:}
\[
\boxed{R_T \leq 8(1 + \ln T) = O(\log T)}
\]

The regret bound is \textbf{logarithmic in $T$} and \textbf{independent of the dimension $d$}. This remarkable dimension-free property arises because:
\begin{enumerate}
    \item The FTL prediction (the running average) has a simple closed form.
    \item The squared Euclidean loss decomposes coordinate-wise, but the key quantity $\|x_t - y_t\|_2^2$ is bounded by a constant (4) independent of $d$, due to the constraint $\|y_t\|_2 \leq 1$.
\end{enumerate}

\end{proof}
        \item[6.]
        \textbf{Problem:} Assuming $\|g_t\|_2 \leq L$ and $\|x\|_2 \leq D, \forall x \in V$. Show that FTRL with regularization $\psi(w) = \frac{1}{2\eta}\|x\|_2^2$ for linear losses $\ell_t(x) = g_t^\top x$, $x \in V$ achieves a sublinear regret.

\bigskip

\textbf{Solution:}

\begin{proof}

\textbf{Step 1: Define the FTRL Algorithm}

The Follow-The-Regularized-Leader (FTRL) algorithm with regularizer $\psi(x) = \frac{1}{2\eta}\|x\|_2^2$ selects:
\[
x_{t+1} = \arg\min_{x \in V} \left\{\sum_{s=1}^{t} \ell_s(x) + \psi(x)\right\} = \arg\min_{x \in V} \left\{\sum_{s=1}^{t} g_s^\top x + \frac{1}{2\eta}\|x\|_2^2\right\}.
\]

\bigskip

\textbf{Step 2: FTRL Regret Lemma}

For FTRL, the regret against any comparator $u \in V$ satisfies:
\[
R_T = \sum_{t=1}^{T}\ell_t(x_t) - \sum_{t=1}^{T}\ell_t(u) \leq \psi(u) - \psi(x_1) + \sum_{t=1}^{T}\left[\ell_t(x_t) - \ell_t(x_{t+1})\right].
\]

Since $\psi(x) = \frac{1}{2\eta}\|x\|_2^2 \geq 0$ and we can initialize $x_1 = 0$ (so $\psi(x_1) = 0$):
\[
R_T \leq \psi(u) + \sum_{t=1}^{T}\left[\ell_t(x_t) - \ell_t(x_{t+1})\right].
\]

\bigskip

\textbf{Step 3: Compute the FTRL Update}

For unconstrained optimization (or when the constraint is not active), taking the gradient:
\[
\nabla_x \left[\sum_{s=1}^{t} g_s^\top x + \frac{1}{2\eta}\|x\|_2^2\right] = \sum_{s=1}^{t} g_s + \frac{1}{\eta}x = 0.
\]

This gives:
\[
x_{t+1} = -\eta \sum_{s=1}^{t} g_s.
\]

If projection onto $V$ is needed, then $x_{t+1} = \Pi_V\left(-\eta \sum_{s=1}^{t} g_s\right)$.

\bigskip

\textbf{Recall (Cauchy-Schwarz Inequality):} For any vectors $a, b \in \mathbb{R}^d$:
$$|a^\top b| \leq \|a\|_2 \cdot \|b\|_2$$

\textbf{Step 4: Bound the Stability Term $\ell_t(x_t) - \ell_t(x_{t+1})$}

For linear losses:
\[
\ell_t(x_t) - \ell_t(x_{t+1}) = g_t^\top x_t - g_t^\top x_{t+1} = g_t^\top(x_t - x_{t+1}).
\]

By Cauchy-Schwarz inequality:
\[
g_t^\top(x_t - x_{t+1}) \leq \|g_t\|_2 \cdot \|x_t - x_{t+1}\|_2 \leq L\|x_t - x_{t+1}\|_2.
\]

\bigskip

\textbf{Step 5: Bound $\|x_t - x_{t+1}\|_2$}

From the FTRL update (unconstrained case):
\begin{align*}
x_t &= -\eta\sum_{s=1}^{t-1}g_s, \\
x_{t+1} &= -\eta\sum_{s=1}^{t}g_s.
\end{align*}

Therefore:
\[
x_{t+1} - x_t = -\eta g_t,
\]
and
\[
\|x_t - x_{t+1}\|_2 = \eta\|g_t\|_2 \leq \eta L.
\]

For the projected case, projection onto a convex set is non-expansive, so the same bound holds.

\bigskip

\textbf{Step 6: Combine the Bounds}

The stability term is bounded by:
\[
\ell_t(x_t) - \ell_t(x_{t+1}) \leq L \cdot \eta L = \eta L^2.
\]

Summing over all rounds:
\[
\sum_{t=1}^{T}\left[\ell_t(x_t) - \ell_t(x_{t+1})\right] \leq \eta L^2 T.
\]

\bigskip

\textbf{Step 7: Total Regret Bound}

For any comparator $u \in V$ with $\|u\|_2 \leq D$:
\[
\psi(u) = \frac{1}{2\eta}\|u\|_2^2 \leq \frac{D^2}{2\eta}.
\]

Therefore, the total regret is:
\[
R_T \leq \frac{D^2}{2\eta} + \eta L^2 T.
\]

\bigskip

\textbf{Step 8: Optimize the Learning Rate $\eta$}

To minimize the bound, we take the derivative with respect to $\eta$ and set it to zero:
\[
\frac{d}{d\eta}\left(\frac{D^2}{2\eta} + \eta L^2 T\right) = -\frac{D^2}{2\eta^2} + L^2 T = 0.
\]

Solving for $\eta$:
\[
\eta^* = \frac{D}{L\sqrt{2T}}.
\]

\bigskip

\textbf{Step 9: Final Regret Bound}

Substituting $\eta^* = \frac{D}{L\sqrt{2T}}$:

\textbf{First term:}
\begin{align*}
\frac{D^2}{2\eta^*} &= \frac{D^2}{2} \cdot \frac{L\sqrt{2T}}{D} = \frac{D^2 \cdot L\sqrt{2T}}{2D} = \frac{DL\sqrt{2T}}{2}
\end{align*}

\textbf{Second term:}
\begin{align*}
\eta^* L^2 T &= \frac{D}{L\sqrt{2T}} \cdot L^2 T = \frac{DLT}{\sqrt{2T}} = \frac{DL\sqrt{T}}{\sqrt{2}} = \frac{DL\sqrt{2T}}{2}
\end{align*}

\textbf{Total:}
\begin{align*}
R_T &\leq \frac{DL\sqrt{2T}}{2} + \frac{DL\sqrt{2T}}{2} = DL\sqrt{2T}
\end{align*}

\bigskip

\textbf{Conclusion:}

\[
\boxed{R_T \leq DL\sqrt{2T} = O\left(\sqrt{T}\right)}
\]

The regret bound $O(\sqrt{T})$ is \textbf{sublinear} in $T$, which means:
\[
\lim_{T \to \infty} \frac{R_T}{T} = \lim_{T \to \infty} \frac{DL\sqrt{2T}}{T} = \lim_{T \to \infty} \frac{DL\sqrt{2}}{\sqrt{T}} = 0.
\]

This guarantees that the average regret $\frac{R_T}{T} \to 0$ as $T \to \infty$, meaning FTRL performs nearly as well as the best fixed decision in hindsight.

\bigskip

\textbf{Summary:}
\begin{itemize}
    \item FTRL update: $x_{t+1} = \Pi_V\left(-\eta\sum_{s=1}^{t}g_s\right)$
    \item Optimal learning rate: $\eta^* = \frac{D}{L\sqrt{2T}}$
    \item Regret bound: $R_T \leq DL\sqrt{2T} = O(\sqrt{T})$ (sublinear)
\end{itemize}

\end{proof}

\end{enumerate}
	
\end{document}
